% Physical pages: 138--144 (printed pages 126--132).

% PCT-SOURCE: pdf=138 print=126 kind=section-title id=sec-ch3-symmetries
\subsection{Symmetries in a Field Theory}
\label{sec:ch3-symmetries}

% PCT-SOURCE: pdf=138 print=126 kind=prose id=p126-symmetries-opening
We now return to the point made at the end of Section \ref{sec:ch1-symmetry-operations}, namely, that if a
unitary operator in $\Hilbert$ is to be interpreted as a symmetry it must have
a reasonable physical meaning.  We first show that the assumption of a
definite transformation law for all the fields of a field theory uniquely
fixes the corresponding transformation on the states up to a phase.

% PCT-SOURCE: pdf=138 print=126 kind=prose id=p126-symmetries-scalar-reduction
For simplicity we prove the theorem only for a hermitian scalar field which
transforms under the parity operator in the usual way.  Simple alterations in
the notation establish similar results in a general field theory for the
operators $P$, $C$, and $T$ given by \eqref{eq:ch1-general-spinor-state-rules}.

% PCT-SOURCE: pdf=139 print=127 kind=theorem id=thm-3-8
\begin{theorem}[3-8]\label{thm:ch3-symmetry-uniqueness}
Let a field theory of a hermitian scalar field $\varphi$ with domain of
definition $D$ be given.  Suppose $U(I_s)$ is a unitary operator which
satisfies

% PCT-SOURCE: pdf=139 print=127 kind=display id=eq-3-58-domain
\begin{equation*}
  U(I_s)D=D
\end{equation*}

% PCT-SOURCE: pdf=139 print=127 kind=display id=eq-3-58
\begin{equation}
  U(I_s)\varphi(x)U(I_s)^{-1}
  =\varphi(x^0,-\mathbf{x})\equiv\varphi(I_sx).
  \tag{3-58}\label{eq:ch3-symmetry-parity-field}
\end{equation}

Then $U(I_s)$ is determined up to a phase, and
$U(I_s)\ket{\Omega}=e^{\ii\alpha}\ket{\Omega}$.  Further,

% PCT-SOURCE: pdf=139 print=127 kind=display id=eq-3-59
\begin{equation}
  U(I_s)U(a,\Lambda)U(I_s)^{-1}
  =U(I_sa,I_s^{-1}\Lambda I_s),
  \tag{3-59}\label{eq:ch3-symmetry-group-conjugation}
\end{equation}
where $U(a,\Lambda)$ is the representation of
$\mathcal{P}_+^\uparrow$.
\end{theorem}

% PCT-SOURCE: pdf=139 print=127 kind=proof id=proof-3-8
\begin{proof}
Choose $\ket{\Psi},\ket{\Phi}\in D$ and consider

% PCT-SOURCE: pdf=139 print=127 kind=display id=eq-3-8-proof-start
\begin{equation*}
\bra{\Psi}\,\varphi(x)U(I_s)U(a,\Lambda)U(I_s)^{-1}
U(I_sa,I_s^{-1}\Lambda I_s)^{-1}\ket{\Phi}.
\end{equation*}

% PCT-SOURCE: pdf=139 print=127 kind=prose id=p127-proof-conjugation
By using \eqref{eq:ch3-symmetry-parity-field} and \eqref{eq:ch3-field-transform}, moving $\varphi$ to the right, we transform this
expression to

% PCT-SOURCE: pdf=139 print=127 kind=display id=eq-3-8-proof-moved-field
\begin{equation*}
\bra{\Psi}\,U(I_s)U(a,\Lambda)U(I_s)^{-1}
U(I_sa,I_s^{-1}\Lambda I_s)^{-1}\varphi(x)\ket{\Phi}.
\end{equation*}

% PCT-SOURCE: pdf=139 print=127 kind=prose id=p127-proof-irreducible
Since, anticipating Theorem \ref{thm:ch4-5-smeared-fields}, we may assume $\varphi$ is irreducible,

% PCT-SOURCE: pdf=139 print=127 kind=display id=eq-3-60
\begin{equation}
  U(I_s)U(a,\Lambda)U(I_s)^{-1}U(I_sa,I_s^{-1}\Lambda I_s)^{-1}
  =\omega\mathbf{1},
  \tag{3-60}\label{eq:ch3-symmetry-phase}
\end{equation}
where $|\omega|=1$.  Applying both sides of \eqref{eq:ch3-symmetry-phase} to the vacuum and using
$U(a,\Lambda)\ket{\Omega}=\ket{\Omega}$ we see that

% PCT-SOURCE: pdf=139 print=127 kind=display id=eq-3-61
\begin{equation}
  U(I_sa,I_s^{-1}\Lambda I_s)U(I_s)\ket{\Omega}
  =\omega^{-1}U(I_s)\ket{\Omega}
  \tag{3-61}\label{eq:ch3-symmetry-vacuum-phase}
\end{equation}
for all $a$, $\Lambda$.  Since the vacuum $\ket{\Omega}$ is the only state
invariant under all $U(a,\Lambda)$,

% PCT-SOURCE: pdf=139 print=127 kind=display id=eq-3-8-vacuum-eigenvalue
\begin{equation*}
  U(I_s)\ket{\Omega}=e^{\ii\alpha}\ket{\Omega}.
\end{equation*}

% PCT-SOURCE: pdf=139 print=127 kind=prose id=p127-proof-phase
Substituting in \eqref{eq:ch3-symmetry-vacuum-phase} we see $\omega=1$, proving \eqref{eq:ch3-symmetry-group-conjugation}.

% PCT-SOURCE: pdf=139 print=127 kind=prose id=p127-proof-extension
To show finally that $U(I_s)$ is determined up to a phase, suppose
$U(I_s)\ket{\Omega}=\ket{\Omega}$.  Then, for states of the form
$\ket{\Psi}=\varphi(f_1)\cdots\varphi(f_n)\ket{\Omega}$, we get

% PCT-SOURCE: pdf=139 print=127 kind=display id=eq-3-8-proof-extension
\begin{equation*}
\begin{aligned}
  U(I_s)\ket{\Psi}
  &={ }U(I_s)\varphi(f_1)U(I_s)^{-1}U(I_s)\varphi(f_2)U(I_s)^{-1}
  \cdots U(I_s)\varphi(f_n)U(I_s)^{-1}U(I_s)\ket{\Omega} \\
  &={ }\varphi(\widehat f_1)\cdots\varphi(\widehat f_n)\ket{\Omega}
  \in D,
\end{aligned}
\end{equation*}
where $\widehat f_j(x)=f_j(I_sx)$.  Thus the operator $U(I_s)$ can be
extended uniquely by linearity and continuity to the whole of $\Hilbert$, and,
so extended, continues to satisfy \eqref{eq:ch3-symmetry-group-conjugation}.
\end{proof}

% PCT-SOURCE: pdf=140 print=128 kind=prose id=p128-discrete-inversions
In the same way the requirements \eqref{eq:ch1-general-spinor-state-rules} for $U(C)$ and $U(I_t)$ lead to the
properties $U(C)\ket{\Omega}=e^{\ii\alpha_C}\ket{\Omega}$,
$U(I_t)\ket{\Omega}=e^{\ii\alpha_t}\ket{\Omega}$, and the ``group''
properties

% PCT-SOURCE: pdf=140 print=128 kind=display id=eq-3-62
\begin{equation}
  U(I_t)U(a,\Lambda)U(I_t)^{-1}
  =U(I_ta,I_t^{-1}\Lambda I_t).
  \tag{3-62}\label{eq:ch3-time-inversion-group}
\end{equation}

% PCT-SOURCE: pdf=140 print=128 kind=display id=eq-3-63
\begin{equation}
  U(C)U(a,\Lambda)U(C)^{-1}=U(a,\Lambda).
  \tag{3-63}\label{eq:ch3-charge-conjugation-group}
\end{equation}

% PCT-SOURCE: pdf=140 print=128 kind=prose id=p128-inversion-properties
Equations \eqref{eq:ch3-symmetry-group-conjugation}, \eqref{eq:ch3-time-inversion-group}, and \eqref{eq:ch3-charge-conjugation-group} and the invariance of
$\ket{\Omega}$ are desirable properties of the operators $U(I_s)$,
$U(I_t)$, and $U(C)$ for them to be interpreted as $P$, $T$, and $C$,
respectively.

% PCT-SOURCE: pdf=140 print=128 kind=prose id=p128-vacuum-phase-convention
It is a standard convention to choose the arbitrary phase factor in the
definition of $U(I_s)$, $U(I_t)$, or $U(C)$ so that

% PCT-SOURCE: pdf=140 print=128 kind=display id=eq-3-64
\begin{equation}
  U(I_s)\ket{\Omega}=\ket{\Omega},\qquad
  U(I_t)\ket{\Omega}=\ket{\Omega},\qquad
  U(C)\ket{\Omega}=\ket{\Omega}.
  \tag{3-64}\label{eq:ch3-inversion-vacuum-convention}
\end{equation}

% PCT-SOURCE: pdf=140 print=128 kind=prose id=p128-symmetry-uniqueness-summary
Then, the results of Theorem \ref{thm:ch3-symmetry-uniqueness} can be summarized: \emph{the requirement
that the fields of a field theory have definite transformation laws under
$U(I_s)$, $U(I_t)$, or $U(C)$ uniquely fixes those operators}.

% PCT-SOURCE: pdf=140 print=128 kind=prose id=p128-symmetry-physical-content
Thus, in a field theory the problem of determining whether a given symmetry
is a reasonable expression of the operations $P$, $T$, or $C$ is reduced to
the problem of understanding the physical content of transformation laws of
the fields under those operations.  This is a somewhat complicated subject
whose full ramifications cannot be completely explored here.  We shall
content ourselves with a series of remarks.

% PCT-SOURCE: pdf=140 print=128 kind=prose id=p128-observable-transformations
For an observable field a transformation law under $P$, $C$, or $T$ is a
direct statement about observable quantities, and no possibility of ambiguity
arises.  For example, in the theory of neutral $\pi$-mesons the transformation
law $\varphi(x)\mapsto-\varphi(I_sx)$ appropriate to pseudo-scalar
$\pi^0$-mesons is physically distinct from
$\varphi(x)\mapsto\varphi(I_sx)$ which would describe scalar $\pi^0$'s.  For
unobservable fields on the other hand, the physical consequences of the
assumption of a given transformation law are only indirectly ascertainable;
an analysis of the super-selection rules of the theory is in general
indispensable for deciding whether the assumption of a different
transformation law actually leads to a physically distinct theory.  For
example, the choices of phase in the definitions \eqref{eq:ch1-dirac-pct-rules} of $P$, $C$, and $T$
for a single spin $\tfrac12$ field $\psi$ are rather arbitrary, as is the
generalization to higher spin; the choice $\psi(x)\mapsto\gamma^0\psi(I_sx)$
is just as good a definition of the parity operation as
$\psi(x)\mapsto-\gamma^0\psi(I_sx)$.  If $U(I_s)$ is the transformation which
produces the first choice, then $RU(I_s)$ will produce the second, if $R$ is
the operator which rotates through an angle $2\pi$ around some axis, i.e., $R$
is $1$ on the states of integer spin and $-1$ on the states of half-odd
integer spin.  In general, two

% PCT-SOURCE: pdf=141 print=129 kind=prose-continuation id=p129-symmetry-operator-equivalence
different choices, $U$ and $U'$, for a symmetry operator are physically
identical if there exists an operator $V$ such that $U=VU'$ and
$\left|\braket{V\Psi}{\Phi}\right|^2=\left|\braket{\Psi}{\Phi}\right|^2$
for all physically realizable states $\ket{\Phi}$, $\ket{\Psi}$.

% PCT-SOURCE: pdf=141 print=129 kind=prose id=p129-anticommuting-fields
A simple example in which the phases of the transformation law under
$U(I_s)$ do have physical consequences is the case of two anti-commuting spin
one-half fields, $\psi_1$ and $\psi_2$.  Suppose
$\psi_1(x)\mapsto\epsilon_1\gamma^0\psi_1(I_sx)$,
$\psi_2(x)\mapsto\epsilon_2\gamma^0\psi_2(I_sx)$ under the operation of
space-inversion.  The transformation $R$ described before can be used to
obtain new fields $\psi'_1$, $\psi'_2$ which have $(\epsilon_1,\epsilon_2)$
replaced by $(-\epsilon_1,-\epsilon_2)$, but the theories with
$(\epsilon_1,\epsilon_2)=(1,+1),(1,-1),(1,\ii)$ are physically distinct.
For example, the quantity
$\psi_1^\dagger(x)\psi_2(y)+\psi_2^\dagger(y)\psi_1(x)$ is a candidate for
an observable in the first two cases, being scalar in the first and
pseudo-scalar in the second, but in the third, it has no well-defined
transformation law.  The third case is interesting because in it
$U(I_s)^2$ commutes with $\psi_2(x)$ and anti-commutes with $\psi_1(x)$.  This
implies the existence of a super-selection rule separating states of the form
$P(\psi_1(f),\psi_2(g),\ldots)\ket{\Omega}$, with $P$ a polynomial odd in
$\psi_1$, which span a subspace $\Hilbert_1$, from those in $\Hilbert_2$ which
is spanned by states with $P$ even in $\psi_1$.  To see this recall that a
physically realizable state $\ket{\Psi}$ must be unchanged when twice
transformed by the parity operator, so if $\ket{\Psi}$ is a vector in a ray
$\Psi$, so is $U(I_s)^2\ket{\Psi}$.  Now $U(I_s)^2$ is $-1$ on $\Hilbert_1$
and $+1$ on $\Hilbert_2$, so a state represented by a vector
$\alpha\ket{\Psi_1}+\beta\ket{\Psi_2}$ with $\alpha\beta\ne0$,
$\ket{\Psi_1}\in\Hilbert_1$, and $\ket{\Psi_2}\in\Hilbert_2$ cannot be
physically realizable.  More generally, if $U(I_s)$, $U(C)$, or $U(I_t)$ are
to be interpreted as $P$, $C$, or $T$, then each of $U(I_s)^2$,
$U(I_t)^2$, $U(C)^2$, $U(I_s)U(C)U(I_s)U(C)$, etc., must be a constant
multiple of the identity on every coherent subspace.  (See Section \ref{sec:ch1-super-selection-rules} for
the definition of physically realizable states and coherent subspaces.)
What this result means in practice is that for certain choices of phases in
$U(I_s)$, $U(I_t)$, or $U(C)$ theories are more general in the sense that
they do not have super-selection rules forced on them by their invariance
under $P$, $C$, or $T$.  It was this kind of special choice which was made in
Section \ref{sec:1.3}.

% PCT-SOURCE: pdf=141 print=129 kind=prose id=p129-multiplicative-symmetry
A common situation in which different phases in the definition of inversions
are physically equivalent is that in which a theory possesses a
\emph{multiplicative symmetry}.  That means that there exists a unitary
operator $V$ such that for some set of complex numbers of absolute value $1$

% PCT-SOURCE: pdf=141 print=129 kind=display id=eq-3-65
\begin{equation}
  V\psi_j(x)V^{-1}=\lambda_j\psi_j(x).
  \tag{3-65}\label{eq:ch3-multiplicative-symmetry}
\end{equation}

% PCT-SOURCE: pdf=141 print=129 kind=prose-continuation id=p129-multiplicative-symmetry-example
The $R$ described previously is an example of a multiplicative symmetry; in
that case $\lambda_j=+1$ for integer spin fields and $-1$ for half-odd
integer spin fields.  By an argument just like that used in the proof of
Theorem \ref{thm:ch3-symmetry-uniqueness} one concludes that $V$ commutes with $U(a,\Lambda)$ and leaves
the vacuum state invariant.  Thus

% PCT-SOURCE: pdf=142 print=130 kind=display id=eq-3-65-selection-rule
\begin{equation*}
\begin{aligned}
  \bra{\Omega}\psi_1(x_1)\cdots\psi_n(x_n)\ket{\Omega}
  &={ }\bra{V\Omega}V\psi_1(x_1)\cdots\psi_n(x_n)\ket{\Omega} \\
  &={ }\lambda_1^{k_1}\lambda_2^{k_2}\cdots
  \bra{\Omega}\psi_1(x_1)\cdots\psi_n(x_n)\ket{\Omega},
\end{aligned}
\end{equation*}
where $k_j$ is the number of times $\psi_j$ appears in the vacuum
expectation value.  Clearly, $\lambda_1^{k_1}\lambda_2^{k_2}\cdots\ne1$
implies that the corresponding vacuum expectation value is zero.  The
vanishing of such vacuum expectation values means that the theory possesses
special selection rules.  In such a theory the choice of $U(I_s)$ to
represent space inversion can be physically equivalent to the choice
$VU(I_s)$.

% PCT-SOURCE: pdf=142 print=130 kind=prose id=p130-collision-theory
There is still another aspect of the transformation law of fields under
symmetries which is essential if their physical interpretation is to be
complete; that is the connection with collision theory.  A straightforward
application of the Haag--Ruelle theory shows that the transformation laws
under $P$, $C$, or $T$ for the fields determine corresponding transformation
laws for the in- and out-fields.  For example, if there is a scalar particle
in the theory, then the associated asymptotic fields satisfy

% PCT-SOURCE: pdf=142 print=130 kind=display id=eq-3-65-in
\begin{equation*}
  U(I_s)\varphi^{\mathrm{in}}(x)U(I_s)^{-1}
  =\varphi^{\mathrm{in}}(I_sx)
\end{equation*}

% PCT-SOURCE: pdf=142 print=130 kind=display id=eq-3-65-c
\begin{equation*}
  U(C)\varphi^{\mathrm{in}}(x)U(C)^{-1}
  =\varphi^{\mathrm{in}\,*}(x)
\end{equation*}

% PCT-SOURCE: pdf=142 print=130 kind=display id=eq-3-65-t
\begin{equation*}
  U(I_t)\varphi^{\mathrm{in}}(x)U(I_t)^{-1}
  =\varphi^{\mathrm{out}}(I_tx)
\end{equation*}

% PCT-SOURCE: pdf=142 print=130 kind=prose id=p130-collision-theory-continuation
and also the same equations with $\varphi^{\mathrm{in}}$ interchanged with
$\varphi^{\mathrm{out}}$.  If in the theory there is a particle with higher
spin, its covariance properties are determined by those of the monomial in
the basic fields in the theory which has a non-zero matrix element between the
vacuum and the state consisting of one such particle.  The operator $U(I_s)$
in \eqref{eq:ch1-dirac-pct-rules} was chosen so that it reverses the momenta of the particles in the
in-states; the same operator has the same effect on the momenta of the
particles in the out-states, as is required physically for the parity operator
[that is, it is independent of time; cf. the remark after \eqref{eq:ch1-space-inversion}].  In the
same way $U(C)$ takes particles to their anti-particles.  It follows that if
$U(I_s)$, $U(C)$ define symmetries they commute with the $S$-operator, giving
rise to a \emph{conservation law}: if $\ket{\Psi^{\mathrm{in}}}$ is an
eigenstate of $U(I_s)$, then so is
$\ket{\Psi^{\mathrm{out}}}=S\ket{\Psi^{\mathrm{in}}}$, and with the same
eigenvalue.  Invariance under $I_t$ is more complicated, since $U(I_t)$
takes an in-state to an out-state with reversed momenta and spins; it leads
to reality conditions on the $S$-matrix elements.

% PCT-SOURCE: pdf=142 print=130 kind=prose id=p130-pct-reconstruction-intro
We now extend the proof of the reconstruction theorem (Theorem \ref{thm:ch3-reconstruction}) to
theories with a discrete symmetry which gives rise to equations

% PCT-SOURCE: pdf=143 print=131 kind=prose-continuation id=p131-pct-reconstruction-intro
like \eqref{eq:pct-3-38} and \eqref{eq:pct-3-39}.  The construction of the corresponding $U(C)$ or
$\Theta$ follows the same lines as the construction of $U(a,\Lambda)$ in
Theorem \ref{thm:ch3-reconstruction}.  We content ourselves with a discussion of the PCT operator.  If
we are given further identities among the $\Wightman$'s, for $P$, $C$, and
$T$ separately, such as \eqref{eq:pct-3-38} for $C$, we can prove the existence of the
corresponding operator in the same way.

% PCT-SOURCE: pdf=143 print=131 kind=theorem id=thm-3-9
\begin{theorem}[3-9]\label{thm:ch3-pct}
Consider a field theory of fields
$\varphi_{(\alpha)(\dot\beta)},\ldots,\psi_{(\mu)(\dot\nu)}$, where
$(\alpha),\ldots,(\mu)$ are collections of undotted indices and
$(\dot\beta),\ldots,(\dot\nu)$ are collections of dotted indices.  Suppose
\eqref{eq:pct-3-39} holds for all vacuum expectation values of the fields, that is,

% PCT-SOURCE: pdf=143 print=131 kind=display id=eq-3-66
\begin{equation}
\begin{aligned}
&\bra{\Omega}\varphi_{(\alpha)(\dot\beta)}(x_1)\cdots
\psi_{(\mu)(\dot\nu)}(x_n)\ket{\Omega} \\
&\quad={ }\ii^F(-1)^J
\overline{\bra{\Omega}\varphi^*_{(\alpha)(\dot\beta)}(-x_1)\cdots
\psi^*_{(\mu)(\dot\nu)}(-x_n)\ket{\Omega}} .
\end{aligned}
\tag{3-66}\label{eq:ch3-pct-vacuum-identity}
\end{equation}
where $F$ is the number of half-odd integer spin fields and $J$ is the total
number of undotted indices in $(\alpha),\ldots,(\nu)$.  Then there exists a
unique (up to a factor) anti-unitary operator $\Theta$ in $\Hilbert$ such that
for any field in the theory

% PCT-SOURCE: pdf=143 print=131 kind=display id=eq-3-67
\begin{equation}
  \Theta^{-1}\varphi_{(\alpha)(\dot\beta)}(f)\Theta
  =(-1)^j\ii^{F(\varphi)}
  \varphi^*_{(\alpha)(\dot\beta)}(\widehat{\overline f}),
  \tag{3-67}\label{eq:ch3-pct-field-transformation}
\end{equation}
where

% PCT-SOURCE: pdf=143 print=131 kind=display id=eq-3-67-test-function
\begin{equation*}
  \widehat{\overline f}(x)=\overline f(-x),\qquad
  (\alpha)=(\alpha_1,\ldots,\alpha_j),\qquad
  (\dot\beta)=(\dot\beta_1,\ldots,\dot\beta_k),
\end{equation*}
and

% PCT-SOURCE: pdf=143 print=131 kind=display id=eq-3-67-parity
\begin{equation*}
  \begin{aligned}
    F^{(\varphi)}&=0 &&\text{if $j+k$ is even},\\
    &=1 &&\text{if $j+k$ is odd}.
  \end{aligned}
\end{equation*}
\end{theorem}

% PCT-SOURCE: pdf=143 print=131 kind=proof id=proof-3-9
\begin{proof}
We can define

% PCT-SOURCE: pdf=143 print=131 kind=display id=eq-3-9-proof-definition
\begin{equation*}
\begin{aligned}
&\Theta\varphi_{(\alpha)(\dot\beta)}(f_1)\cdots
\psi_{(\mu)(\dot\nu)}(f_n)\ket{\Omega} \\
&\quad={ }\Theta\varphi_{(\alpha)(\dot\beta)}(f_1)\Theta^{-1}\Theta\cdots
\Theta\psi_{(\mu)(\dot\nu)}(f_n)\Theta^{-1}\Theta\ket{\Omega} \\
&\quad={ }(-\ii)^F(-1)^J
\varphi^*_{(\alpha)(\dot\beta)}(\widehat{\overline{f_1}})\cdots
\psi^*_{(\mu)(\dot\nu)}(\widehat{\overline{f_n}})\ket{\Omega}\in H.
\end{aligned}
\end{equation*}

% PCT-SOURCE: pdf=143 print=131 kind=prose id=p131-pct-antilinearity
This equation can be extended by anti-linearity to the whole of the space $H$
(defined in Section \ref{sec:ch3-reconstruction}).  So defined, $\Theta$ is anti-unitary, by virtue of
\eqref{eq:ch3-pct-vacuum-identity}, which leads to

% PCT-SOURCE: pdf=143 print=131 kind=display id=eq-3-9-antiunitary
\begin{equation*}
  \braket{\Theta f}{\Theta g}=\overline{\braket{f}{g}}
\end{equation*}
for all finite sequences $f=(f_0,f_1,\ldots)$, $g=(g_0,g_1,\ldots)$ of test
functions.  In fact $\Theta$ defines a mapping of the quotient space $H/H_0$
obtained by identifying vectors whose difference has zero norm in $H$ (see
Section \ref{sec:ch3-reconstruction}).  For suppose $f\sim g$ in the sense that
$\lVert f-g\rVert=0$.  This means that a certain sum of vacuum expectation
values smeared with the $f$'s and $g$'s vanishes.

% PCT-SOURCE: pdf=144 print=132 kind=proof-continuation id=proof-3-9-continuation
We may replace these vacuum expectation values by the PCT transformed ones,
using \eqref{eq:ch3-pct-vacuum-identity}.  The new equation is then a statement that
$\lVert\Theta f-\Theta g\rVert=0$, i.e., $\Theta f\sim\Theta g$.
Therefore $\Theta$ defines an anti-unitary operator on $H/H_0$, which may be
regarded as a subset of $\Hilbert$.  The extension to all of $\Hilbert$ is done
by continuity, and, so defined, is anti-unitary.
\end{proof}

% PCT-SOURCE: pdf=144 print=132 kind=prose id=p132-pct-conclusion
Theorem \ref{thm:ch3-pct} shows that the presence of PCT symmetry in a field theory is
equivalent to the validity of the identities \eqref{eq:ch3-pct-vacuum-identity}.  It is shown in the next
chapter that the identities \eqref{eq:ch3-pct-vacuum-identity} hold in every local field theory.  This is
obviously a very important result.
